\rhead{RB201: Mechanics of Robotics}
\phantomsection 
\section*{Mechanics of Robotics (RB201)}
\label{major:RB_1}

\noindent \small{\hyperref[main:scheme]{$\leftarrow$ to Scheme}} \vspace{0.5cm}

\noindent \textbf{Credits:} 3 (2-1-0)

\subsection*{Course Content}
\noindent\textbf{Unit 1} \\
\textit{Spatial Descriptions and Transformations:} Introduction to robotics, Robot components (links, joints), and Degrees of Freedom (DOF). Position and Orientation representation in 3D space. Coordinate Frames. Rotation Matrices and their properties. Homogeneous Transformation Matrices for representing combined rotation and translation. Composition of transformations.

\vspace{0.3cm}

\noindent\textbf{Unit 2} \\
\textit{Forward and Inverse Kinematics:} Kinematic chains. Forward Kinematics: Calculating the end-effector's position and orientation from given joint angles. Denavit-Hartenberg (D-H) convention for systematically assigning coordinate frames to manipulators. Inverse Kinematics: Calculating the required joint angles to achieve a desired end-effector pose. Analytical vs. Numerical solutions. Workspace analysis.

\vspace{0.3cm}

\noindent\textbf{Unit 3} \\
\textit{Velocity Kinematics - The Jacobian:} Linear and angular velocity of rigid bodies. The Jacobian matrix relating joint velocities to the end-effector's linear and angular velocities. Derivation of the Jacobian. Singularities: Configurations where the manipulator loses degrees of freedom and the Jacobian is not invertible. Relationship between Jacobian and static forces.

\vspace{0.3cm}

\noindent\textbf{Unit 4} \\
\textit{Manipulator Dynamics:} Introduction to dynamics. Newton-Euler formulation for a single rigid body. Lagrangian Dynamics: A systematic approach based on kinetic and potential energy. Derivation of the equations of motion for simple manipulators. The structure of the dynamic model: Inertia matrix, Coriolis and Centrifugal terms, and Gravity vector.

\vspace{0.3cm}

\noindent\textbf{Unit 5} \\
\textit{Trajectory Planning and Introduction to Control:} Path vs. Trajectory. Trajectory generation in joint-space and Cartesian-space. Polynomial trajectories (cubic, quintic) for smooth point-to-point motion. Introduction to robot control concepts. The need for feedback control. Proportional-Derivative (PD) control for joint-space motion. Introduction to different actuation systems (electric, hydraulic, pneumatic).

\newpage
\rhead{Curriculum Schema}